\section{Introduction and Preliminaries}

It is sometimes said that model predictive control (MPC) is the proportional--integral--derivative (PID) control of the 21st century.
While it is true that MPC has gained a lot of attention in the last decades, we still believe that PID control remains the workhorse of the 21st century, simply because it is so easy to implement.
Implementing an MPC is not.
This paper aims to demystify the implementation of MPC and to shed light on its many wonderful mysteries.

\subsection{Linear Systems}

Most physical systems can be described in continuous time by an ordinary differential equation
\[
\dot{x}(t) = f\big(x(t), u(t)\big),
\]
with the state $x(t) \in \R{n_x}$ and the input $u(t) \in \R{n_u}$.
For ease of notation, the time dependence of the states $x$ and inputs $u$ will be dropped from here on.

\paragraph{LTI systems.}
A linear time-invariant (LTI) system is described by
\[
\dot{x} = A_c x + B_c u,
\]
with constant matrices $A_c \in \R{n_x \times n_x}$ and $B_c \in \R{n_x \times n_u}$, where the subscript $c$ denotes continuous time.
Its discrete-time counterpart, derived in the next subsection, reads
\[
x_{k+1} = A_d x_k + B_d u_k,
\]
with constant matrices $A_d \in \R{n_x \times n_x}$ and $B_d \in \R{n_x \times n_u}$, where the subscript $d$ denotes discrete time.
Linear systems obey the superposition principle.
LTI systems additionally have a closed-form solution, and their properties such as stability, controllability and observability can be checked with simple matrix computations.

\paragraph{LTV systems.}
In a linear time-varying (LTV) system, the matrices are allowed to change over time
\[
\dot{x} = A_c(t) x + B_c(t) u + b_c(t),
\]
where we additionally allow an affine term $b_c(t) \in \R{n_x}$.
Strictly speaking, such a system is affine, but we follow common usage and still call it LTV.
In discrete time, this becomes
\[
x_{k+1} = A_k x_k + B_k u_k + b_k,
\]
with stage-dependent $A_k$, $B_k$ and $b_k$.

\paragraph{Linearization along a trajectory.}
The most important source of LTV systems is the linearization of a nonlinear system around a nominal trajectory.
Let $(x^\circ, u^\circ)$ be such a nominal trajectory, for instance a reference trajectory.
A first-order Taylor expansion of the vector field $f$ around this trajectory yields
\[
\dot{x} \approx f\big(x^\circ, u^\circ\big) + \left. \frac{\partial f}{\partial x} \right|_{(x^\circ(t), u^\circ(t))} \big(x - x^\circ\big) + \left. \frac{\partial f}{\partial u} \right|_{(x^\circ(t), u^\circ(t))} \big(u - u^\circ\big),
\]
with the Jacobians
\[
A_c(t) = \left. \frac{\partial f}{\partial x} \right|_{(x^\circ(t), u^\circ(t))},
\qquad
B_c(t) = \left. \frac{\partial f}{\partial u} \right|_{(x^\circ(t), u^\circ(t))}.
\]
Since the Jacobians are evaluated at different points along the trajectory, they vary with time, and the linearized model is an LTV system, unless $(x^\circ, u^\circ)$ is an equilibrium, in which case the linearization is LTI.
Collecting the terms independent of $x$ and $u$ gives the affine form from above with the offset
\[
b_c(t) = f\big(x^\circ, u^\circ\big) - A_c(t) x^\circ - B_c(t) u^\circ.
\]
Alternatively, the model can be written in deviation coordinates $\Delta x = x - x^\circ$, $\Delta u = u - u^\circ$, which gives
\[
\Delta \dot{x} \approx A_c(t)\, \Delta x + B_c(t)\, \Delta u + \big( f(x^\circ, u^\circ) - \dot{x}^\circ \big).
\]
If the nominal trajectory satisfies the nonlinear dynamics, i.e.\ $\dot{x}^\circ = f(x^\circ, u^\circ)$, the bracketed offset vanishes.
Otherwise it measures the mismatch of the nominal trajectory.
In either representation, the linear model is only a good approximation close to the nominal trajectory.
Repeatedly linearizing around the latest solution and solving the resulting linear-quadratic optimal control problem (OCP) is the basis of sequential quadratic programming (SQP) and real-time iteration schemes for nonlinear MPC~\cite{diehl2005real}, which is why efficient solvers for LTV problems are relevant far beyond linear plants.

\subsection{Discretization}

Controllers are implemented or simulated on digital hardware, where inputs are computed and applied at discrete sampling instants.
We therefore require a discrete-time model that maps the state and input at stage $k$ to the state at stage $k+1$.
We will start from a given continuous-time LTI system and arrive at a discrete-time equivalent, which, for inputs held constant between samples, reproduces the sampled continuous-time state exactly.

\paragraph{The transition matrix.}
Without any inputs the dynamics reduce to the autonomous system
\[
  \dot{x} = A_c x, \qquad x(0)=x_0,
\]
with the well-known solution
\[
  x(t) = \Phi(t)x_0 ,
\]
where
\[
  \Phi(t) = \sum_{i=0}^{\infty} A_c^i \frac{t^i}{i!} = e^{A_c t}
\]
is the so-called transition matrix.
Its $i$-th column is the free response to the initial state $e_i$, and by linearity this gives the solution above for any $x_0$.
Differentiating the series term by term shows $\dot{\Phi}(t) = A_c \Phi(t)$.
Furthermore, it has the important properties $\Phi(0)=I$, $\Phi(t)\Phi(\tau)=\Phi(t+\tau)$ and $\Phi^{-1}(t)=\Phi(-t)$.

\paragraph{Solution via variation of constants.}
With inputs, the state no longer follows the free response.
We describe it in the moving basis of free solutions and let the coordinates vary over time,
\[
  x(t) = \Phi(t)\, c(t), \qquad c(0) = x_0.
\]
Differentiating and inserting $\dot{\Phi} = A_c \Phi$ gives
\[
  \dot{x} = \dot{\Phi}\, c + \Phi\, \dot{c} = A_c \Phi c + \Phi \dot{c} = A_c x + \Phi\, \dot{c} .
\]
Comparing with the true dynamics $\dot{x} = A_c x + B_c u$ yields
\[
  \Phi\,\dot{c} = B_c u
\]
and finally
\[
    \dot{c}(t) = \Phi^{-1}(t) B_c\, u(t) = e^{-A_ct} B_c\, u(t) .
\]
Integrating both sides using the fundamental theorem of calculus yields
\[
    c(t) - c(0) = \int_0^{t} \Phi(-\tau) B_c\, u(\tau)\, d\tau .
\]
Using the inverse transformation $c(t)=\Phi^{-1}(t)x(t)=\Phi(-t)x(t)$ and $c(0) = x_0$, we arrive at
\[
    \Phi(-t)x(t) = x_0 + \int_0^{t} \Phi(-\tau) B_c\, u(\tau)\, d\tau .
\]
Multiplying by $\Phi(t)$ and using $\Phi(t)\Phi(-\tau) = \Phi(t-\tau)$, the general solution is given by
\[
    x(t) = \Phi(t)x_0 + \int_0^{t} \Phi(t-\tau) B_c\, u(\tau)\, d\tau .
\]
Since the system is time-invariant, the same holds for any initial time $t_0$, i.e.\ $x(t) = \Phi(t-t_0)x(t_0) + \int_{t_0}^{t} \Phi(t-\tau) B_c\, u(\tau)\, d\tau$.

\paragraph{Zero-order hold.}
Let $T_s > 0$ be the sampling time, $t_k := k T_s$ the sampling instants and $x_k := x(t_k)$.
The input is computed only at the sampling instants and held constant in between, i.e.
\[
    u(t) = u_k \quad \text{for } t\in[t_k, t_{k+1}) .
\]
We can therefore compute the exact transition from the state $x_k$ to $x_{k+1}$ as
\[
    x(t_{k+1}) = \Phi(t_{k+1}-t_k)x(t_k) + \int_{t_k}^{t_{k+1}} \Phi(t_{k+1}-\tau) d\tau\, B_c\, u_k,
\]
and using $\sigma = t_{k+1} - \tau$, with $d\sigma = -d\tau$, where the limits swap from $\tau = t_k, t_{k+1}$ to $\sigma = T_s, 0$ and cancel the sign,
\begin{equation} \label{eq:zoh}
    x_{k+1} = \Phi(T_s)x_k + \int_{0}^{T_s} \Phi(\sigma) d\sigma\, B_c\, u_k .
\end{equation}
Thus we arrive at the zero-order hold (ZOH) discretization, which is exact at the sampling instants $t_k$,
\[
    x_{k+1} = A_dx_k + B_du_k ,
    \qquad
    A_d := \Phi(T_s) = e^{A_c T_s},
    \qquad
    B_d := \int_0^{T_s} \Phi(\sigma)\, d\sigma\, B_c .
\]
Here, $A_d$ is the transition matrix over one sampling period, so $A_d x_k$ is the free response starting from $x_k$.
$B_d$ is the convolution of the impulse response $\Phi(\sigma)B_c$ with a unit-height pulse of width $T_s$, evaluated at the end of the pulse, so $B_d u_k$ is the forced response (from zero initial state) to the input $u_k$.

\paragraph{Computation via Van Loan.}
If $A_c$ is invertible, the discrete-time input matrix can be computed as $B_d = A_c^{-1} \big( e^{A_c T_s} - I \big) B_c$, but this formula fails for singular $A_c$, which is common in practice.
Any system containing an integrator has a singular $A_c$.
A simple trick~\cite{van1978computing} avoids evaluating the integral explicitly.
Under the ZOH assumption, the input is constant on each sampling interval, i.e.\ $\dot{u}(t) = 0$ for $t \in (t_k, t_{k+1})$.
We can therefore treat the input as an additional state and augment the state to $\tilde{x} = [x; u]$, which follows the autonomous LTI dynamics
\[
\dot{\tilde{x}} = M \tilde{x},
\qquad
M :=
\begin{bmatrix}
    A_c & B_c \\
    0   & 0
\end{bmatrix}.
\]
Since the augmented system is autonomous, its solution over one interval is simply $\tilde{x}(t_{k+1}) = e^{M T_s} \tilde{x}(t_k)$.
To see how $A_d$ and $B_d$ appear in this matrix exponential, note that the powers of $M$ have the block structure
\[
M^0 = I,
\qquad
M^i =
\begin{bmatrix}
    A_c^i & A_c^{i-1} B_c \\
    0     & 0
\end{bmatrix},
\quad i \geq 1,
\]
which is easily verified by induction.
Inserting this into the series gives
\[
e^{M T_s} = \sum_{i=0}^{\infty} \frac{T_s^i}{i!} M^i =
\begin{bmatrix}
    \sum_{i=0}^{\infty} \frac{T_s^i}{i!} A_c^i & \sum_{i=1}^{\infty} \frac{T_s^i}{i!} A_c^{i-1} B_c \\
    0 & I
\end{bmatrix}.
\]
The upper left block is $e^{A_c T_s} = A_d$.
For the upper right block, we shift the index and use $\int_0^{T_s} \frac{\sigma^i}{i!} \, d\sigma = \frac{T_s^{i+1}}{(i+1)!}$ to integrate the series term by term
\[
\sum_{i=1}^{\infty} \frac{T_s^i}{i!} A_c^{i-1} B_c
= \sum_{i=0}^{\infty} \frac{T_s^{i+1}}{(i+1)!} A_c^{i} B_c
= \int_0^{T_s} \sum_{i=0}^{\infty} \frac{\sigma^i}{i!} A_c^{i} \, d\sigma \, B_c
= \int_0^{T_s} e^{A_c \sigma} \, d\sigma \, B_c
= B_d.
\]
Hence, both discrete-time matrices are obtained from a single matrix exponential
\begin{equation} \label{eq:zoh_expm}
\exp\left(
\begin{bmatrix}
    A_c & B_c \\
    0   & 0
\end{bmatrix}
T_s
\right)
=
\begin{bmatrix}
    A_d & B_d \\
    0   & I
\end{bmatrix}.
\end{equation}
A constant affine term $b_c$ in the dynamics $\dot{x} = A_c x + B_c u + b_c$ is included in the same way by appending $b_c$ as an additional column, which corresponds to augmenting the state with a constant $1$, i.e.\ $\tilde{x} = [x; u; 1]$,
\[
\exp\left(
\begin{bmatrix}
    A_c & B_c & b_c \\
    0   & 0   & 0   \\
    0   & 0   & 0
\end{bmatrix}
T_s
\right)
=
\begin{bmatrix}
    A_d & B_d & b_d \\
    0   & I   & 0   \\
    0   & 0   & 1
\end{bmatrix},
\qquad
b_d = \int_0^{T_s} e^{A_c \sigma} \, d\sigma \, b_c,
\]
which yields the discrete-time affine system $x_{k+1} = A_d x_k + B_d u_k + b_d$.
For an LTV model, the matrices are typically frozen at $t_k$ on each sampling interval (or the variational equations are integrated numerically), and the augmented exponential is evaluated per stage to obtain $A_k$, $B_k$ and $b_k$.
This discretization is then no longer exact.

\newpage
